% =============================================================================
% SUPPLEMENT NOTE S2 -- Canonical-confusion reciprocity proof
% (R36 s6: "S2 canonical-confusion reciprocity proof, pseudoinverse cases,
%  determinant identities")
% =============================================================================
% R36 CROSSWALK: from R35 old S7 (reciprocity + DLGI estimator), the RECIPROCITY
%   PROOF half moves here; the estimator half moves to R36 S3.
%
% CONTENT MANIFEST (R36 s6 "S2"):
%   * Full proof of the Canonical-Confusion Ledger Reciprocity Theorem.
%   * Pseudoinverse cases (rank-deficient A or B).
%   * Determinant identities.
%
% THEOREM NAME DISCIPLINE (register B3 / R33 s7): theorem =
%   "Canonical-Confusion Ledger Reciprocity Theorem"; method = DLGI. Distinct.
%   Forbidden: "information conservation", "free medium information", additive law.
%
% PORTED WHOLESALE FROM:
%   * R33 Thm 2 (materials map s4b): K=A^{-1/2}CB^{-1/2};
%     J_x=A^{1/2}(I-KK^T)A^{1/2}, J_theta=B^{1/2}(I-K^TK)B^{1/2};
%     det(J_x)/det(A)=prod(1-kappa_i^2)=det(J_theta)/det(B).
%   * results/round63_next/DUAL_LEDGER_PROBE/t3_reciprocity.json
%     (reciprocity det exact to 1e-15; ||K|| schedule-invariant 0.2%).
% =============================================================================
\section{Canonical-confusion reciprocity: proof and determinant identities}
\label{supp:reciprocity}

% TODO: full reciprocity proof; shared nonzero eigenvalues of KK^T and K^TK;
% pseudoinverse (rank-deficient) cases; determinant identities; the O4-A
% "C=0 protects both ledgers" corollary.
